\section{Uniform separation and general diagonalization}
\label{sec:separation}

We now prove Theorem~\ref{thm:individual-to-strong}.  The strict
support bound in Proposition~\ref{prop:shrinking-P} initially
depends on the inputs.  Products turn it into a universal constant
below one; iteration then makes the support compression arbitrarily
small.  A final cutoff comparison gives coefficients whose norms
are bounded independently of that requested smallness.

\begin{lemma}[Products preserve individual proper infiniteness]
\label{lem:product-P}
Let \((C_i)_{i\in I}\) be any family of \(C^*\)-algebras
satisfying \(\propP\).  Their bounded product
\(D=\prod_{i\in I}C_i\) satisfies \(\propP\).
More precisely, for a positive contraction \(d=(d_i)\in D\)
and \(0<\delta<1/2\), there is a row
\(s=(s_i)\in M_{1,2}(D)\) such that
\begin{equation}\label{eq:product-row}
 s^*ds=(d-2\delta)_+\oplus(d-2\delta)_+,\qquad
 \|s\|\le\delta^{-1/2}.
\end{equation}
\end{lemma}

\begin{proof}
For an empty index set, \(D=0\) and the zero row has the asserted
properties.  Otherwise, for each \(i\), put
\[
 h_i=(d_i-\delta)_+,\qquad
 k_i=(d_i-2\delta)_+\oplus(d_i-2\delta)_+.
\]
The positive \(h_i\) is properly infinite.  Apply
Lemma~\ref{lem:transport} to \(h_i\oplus h_i\precsim h_i\)
with cutoff \(\delta\).  If \(v_i\) is the resulting
rectangular partial isometry, then
\[
 r_i=v_i k_i^{1/2}\in M_{1,2}(C_i),\qquad
 r_i^*r_i=k_i,\qquad e_{h_i}r_i=r_i,\qquad \|r_i\|\le1.
\]
The algebra membership follows from transport regularity; no
membership of \(v_i\) itself is needed.  Define
\[
 \phi_\delta(t)=
 \begin{cases}
  t/\delta^{3/2},&0\le t\le\delta,\\
  t^{-1/2},&t\ge\delta.
 \end{cases}
\]
This continuous function vanishes at zero, is bounded by
\(\delta^{-1/2}\), and satisfies
\(t^{1/2}\phi_\delta(t)=1\) for \(t>\delta\).  Consequently
\[
 s_i=\phi_\delta(d_i)r_i\in M_{1,2}(C_i),\qquad
 d_i^{1/2}s_i=r_i,\qquad
 \|s_i\|\le\delta^{-1/2}.
\]
The uniform bound places \((s_i)\) in \(M_{1,2}(D)\).
Coordinate functional calculus gives
\eqref{eq:product-row}.  Its right-hand side converges to
\(d\oplus d\) in norm as \(\delta\downarrow0\), proving
proper infiniteness of \(d\).  Scaling covers every positive
element of \(D\).
\end{proof}

The general product-comparison criterion of
Antoine--Perera--Thiel \cite[Lemma~5.7]{AntoinePereraThiel2020}
also implies this permanence result: for each cutoff
\(\varepsilon\), one can take a target cutoff
\(\delta=\varepsilon/2\) and use proper infiniteness in every
coordinate.  The proof above records the particular bounded
rows that make this inference possible.  Coordinatewise
comparison without a common cutoff bound would not suffice.

\begin{lemma}[A uniform support contraction]\label{lem:uniform-shrinking}
There exists \(\theta\in(0,1)\) such that for every
\(C^*\)-algebra \(C\) satisfying \(\propP\), every pair of positive
contractions \(a,y\in C\), and every contraction \(c\in C\),
there are positive contractions \(p\in\Her(a)\),
\(q\in\Her(y)\) such that, for every representation \(\pi\) of
\(C\),
\[
 \|\pi(a)\|>1/2\ \Longrightarrow\ \|\pi(p)\|=1,
 \qquad
 \|\pi(y)\|>1/2\ \Longrightarrow\ \|\pi(q)\|=1,
\]
and
\begin{equation}\label{eq:uniform-support}
 \|e_p c e_q\|\le\theta.
\end{equation}
\end{lemma}

\begin{proof}
If no such bound existed, for each \(m\ge2\) we could choose
an algebra \(C_m\) satisfying \(\propP\) and contractions
\(a_m,y_m,c_m\), with the first two positive, such that every
pair with the stated membership and simultaneous signals has
support norm greater than \(1-1/m\).  Form their bounded
product \(D\) and the three elements \(a=(a_m)\),
\(y=(y_m)\), \(c=(c_m)\in D\).  Lemma~\ref{lem:product-P}
gives \(\propP\) for \(D\), so
Proposition~\ref{prop:shrinking-P} produces \(p,q\) with the
signals and \(\kappa=\|e_pce_q\|<1\).

For each \(m\), the coordinate homomorphism \(D\to C_m\)
preserves hereditary membership, and its normal extension sends
the two support projections to the coordinate supports.
Thus the coordinate support norm is at most \(\kappa\).
If \(\pi\) is any representation of \(C_m\), its composition
with the coordinate homomorphism is a representation of \(D\),
so both signals hold for the coordinate pair as well.  Choosing
\(m\) with \(1-1/m>\kappa\) contradicts its selection.
The resulting constant can be increased, if necessary, to lie
in \((0,1)\).
\end{proof}

This proof uses the normal coordinate maps, not an identification
of \(D^{**}\) with a product of biduals.  The input \(c\) in
Lemma~\ref{lem:uniform-shrinking} belongs to \(C\).
This ensures that the sequence \((c_m)\) belongs to \(D\);
a sequence of arbitrary individual unitizations need not belong
to \(D\) or its unitization.  We need the uniform statement
only in the stated form.

\begin{proposition}[Arbitrarily small support compression]
\label{prop:iterated-separation}
Suppose \(C\) satisfies \(\propP\).  For positive contractions
\(a_*,b_*\in C\), a contraction \(c_*\in C\), and
\(\zeta>0\), there are positive contractions
\(p\in\Her(a_*)\), \(q\in\Her(b_*)\) such that
\[
 \|e_p c_*e_q\|<\zeta
\]
and, for every representation \(\pi\),
\[
 \|\pi(a_*)\|>1/2\ \Longrightarrow\ \|\pi(p)\|=1,
 \qquad
 \|\pi(b_*)\|>1/2\ \Longrightarrow\ \|\pi(q)\|=1.
\]
\end{proposition}

\begin{proof}
Fix \(\theta\) from Lemma~\ref{lem:uniform-shrinking}.
On \([0,1]\) set
\[
 f(t)=(2t-1)_+,\qquad g(t)=\min\{1,(4t-1)_+\}.
\]
Both vanish near zero, \(f(1)=1\), and \(g=1\) on
\([1/2,1]\).  For a positive contraction \(h\),
\begin{equation}\label{eq:iteration-cutoffs}
 g(h)=e_hg(h),\qquad g(h)e_{f(h)}=e_{f(h)}.
\end{equation}
Apply Lemma~\ref{lem:uniform-shrinking} once to
\(a_*,b_*,c_*\), obtaining \(p_1,q_1\) with support norm
at most \(\theta\).  Suppose inductively that
\(\|e_{p_k}c_*e_{q_k}\|\le\theta^k\).  The next normalized
input is
\begin{equation}\label{eq:iteration-normalization}
 c^{(k)}=\theta^{-k}g(p_k)c_*g(q_k)\in C,\qquad
 \|c^{(k)}\|\le1.
\end{equation}
Apply the uniform lemma to \(f(p_k),f(q_k),c^{(k)}\).
The new supports lie under those of \(f(p_k)\) and \(f(q_k)\);
by \eqref{eq:iteration-cutoffs},
\[
 e_{p_{k+1}}c_*e_{q_{k+1}}
 =\theta^k e_{p_{k+1}}c^{(k)}e_{q_{k+1}},
 \qquad
 \|e_{p_{k+1}}c_*e_{q_{k+1}}\|\le\theta^{k+1}.
\]
The hereditary algebras are nested.  In every representation,
norm one of \(p_k\) implies norm one of \(f(p_k)\), by
spectral mapping and \(f(1)=1\), and similarly for \(q_k\).
Thus the two original signals persist at every step.
All positive contractions needed as first inputs are rich under \(\propP\).
Choose \(k\) with \(\theta^k<\zeta\).
\end{proof}

\begin{proposition}[Comparison with bounded coefficients]
\label{prop:bounded-diagonalization}
Suppose \(C\) satisfies \(\propP\).  Let \(a,b\in C_+\) be
contractions and let \(c\in C\).  Given
\[
 0<\delta<1/2,\qquad 0<\rho\le2\delta,\qquad \zeta>0,
\]
there exist \(s,t\in C\) such that
\begin{align}
 \|s^*as-a\|,\ \|t^*bt-b\|&<2\delta+\rho,
 \label{eq:bounded-diagonals}\\
 \|s^*ct\|&<\zeta,\qquad
 \|s\|,\ \|t\|\le\delta^{-1/2}.
 \label{eq:bounded-cross}
\end{align}
\end{proposition}

\begin{proof}
Let \(a_*,b_*\) be ramps of \(a,b\), zero on
\([0,\delta]\) and one on \([2\delta,1]\).
Put \(M_c=\max\{1,\|c\|\}\).  Choose a finite iteration in
Proposition~\ref{prop:iterated-separation} so that its outputs
satisfy the signals and
\begin{equation}\label{eq:bounded-chosen-separation}
 \|e_p c e_q\|<\delta\zeta.
\end{equation}
For example, apply the construction to \(c/M_c\) and take
\(k\) with \(M_c\theta^k<\delta\zeta\).
This choice precedes the comparison coefficients below.

It remains to realize the diagonal cutoffs through \(p\) and \(q\)
with a coefficient bound depending only on \(\delta\) and \(\rho\).

Put \(d_a=(a-2\delta)_+\).  If an irreducible representation
detects \(d_a\), its norm on \(a\) is greater than \(2\delta\).
It then has norm one on \(a_*\), and the signal gives norm
one on \(p\).  Lemma~\ref{lem:ideals-cutoffs}(ii) yields
\(d_a\in\Id(p)\).  Since
\[
 e_p\le e_{a_*}\le\mathbf1_{(\delta,\infty)}(a),
 \qquad pap\ge\delta p^2,
\]
we have \(\Id(pap)=\Id(p)\): one inclusion is immediate, and
the displayed order gives the other.  The positive \(pap\)
is properly infinite under \(\propP\), so
Lemma~\ref{lem:comparison-rules}(iii) gives
\(d_a\precsim pap\).

Compressing a stabilization comparison coefficient to the scalar
source and target corners supplies \(w\in C\) with
\[
 \|w^*papw-d_a\|<\rho.
\]
Set \(s=pw\).  Then \(e_ps=s\), \(s^*as=w^*papw\), and
\begin{equation}\label{eq:final-coefficient-bound}
 \delta s^*s\le s^*as,\qquad
 \|s\|^2
 \le\frac{\|d_a\|+\rho}{\delta}
 \le\frac{1-2\delta+\rho}{\delta}
 \le\frac1\delta.
\end{equation}
In particular the more general estimate
\(\|pw\|^2\le(1+\rho)/\delta\) holds without the
restriction \(\rho\le2\delta\).  The same argument for \(b,q\)
gives \(t=qw'\), \(e_qt=t\), and the identical bounds.
Because \(\|a-d_a\|\le2\delta\), these comparisons prove
\eqref{eq:bounded-diagonals}.  Finally
\[
 \|s^*ct\|
 \le\|s\|\|t\|\|e_p c e_q\|
 \le\delta^{-1}\|e_p c e_q\|<\zeta.
\]
The bound in \eqref{eq:final-coefficient-bound} depends on
\(\delta,\rho\), not on \(p,q\) or the separation accuracy.
This is why the preceding choice
\eqref{eq:bounded-chosen-separation} remains valid after the
comparison coefficients are found.
\end{proof}

\begin{proof}[Proof of Theorem~\ref{thm:individual-to-strong}]
For arbitrary \(a,b\in C_+\), \(c\in C\), and
\(\varepsilon>0\), put \(M=\max\{1,\|a\|,\|b\|\}\).
Choose \(0<\delta<1/2\) and \(0<\rho\le2\delta\)
with \(M(2\delta+\rho)<\varepsilon\).  Apply
Proposition~\ref{prop:bounded-diagonalization} to \(a/M,b/M,c\)
with \(\zeta=\varepsilon\).  Multiplying the two diagonal
inequalities by \(M\) gives all three inequalities in
\eqref{eq:main-diagonalization}.

For positive \(x,y\in C\), specialize these inequalities to
\((a,b,c)=(x^2,y^2,xy)\).  They are precisely
\eqref{eq:SPI-definition}, so \(C\) is strongly purely infinite.
This uses the unrestricted definition directly.  If one instead
writes a block matrix
\(Z=\left(\begin{smallmatrix}a&c\\c^*&b\end{smallmatrix}\right)\),
the norm of
\(\operatorname{diag}(s,t)^*Z\operatorname{diag}(s,t)-(a\oplus b)\)
is at most the maximum of the two diagonal errors plus
\(\|s^*ct\|\).  Hence the same estimates also give diagonal
approximation of every such block matrix, in particular of every
positive one.  The argument includes zero elements and the zero
algebra.
\end{proof}
